\section{Aplicación: Proyecto}

\subsection{Descripción del alcance del proyecto}
\subsubsection{(i) Alcance general del Proyecto a Implementar}
El proyecto consiste en el diseño e implementación de servicios de consultoría especializados en Sistemas de Información Geográfica (GIS), orientados a mejorar la toma de decisiones mediante análisis espacial, interoperabilidad tecnológica y capacitación de usuarios. Incluye:
\begin{itemize}
  \item Análisis de necesidades y requisitos del cliente.
  \item Desarrollo de soluciones GIS personalizadas.
  \item Integración de plataformas GIS con sistemas existentes.
  \item Capacitación básica a usuarios finales.
  \item Documentación técnica y manuales de usuario.
\end{itemize}

\subsubsection{(ii) Historias de usuario del proyecto}
Tomando como referencia el ejemplo de Scrum del otro documento, podrían definirse historias como:
\begin{description}
  \item[HU001] Como cliente, quiero integrar mi plataforma de gestión con el sistema GIS para visualizar datos geoespaciales en tiempo real.
  \item[HU002] Como usuario final, necesito capacitación básica en el uso de WebGIS para realizar consultas espaciales autónomas.
  \item[HU003] Como administrador, requiero un panel de control con métricas de uso y cumplimiento de actividades obligatorias en la plataforma GIS.
\end{description}

\subsubsection{(iii) Roles Scrum}
	extbf{Product Owner:} Representante del cliente o área solicitante (ej: Dra. Rosmery R.).\\
	extbf{Scrum Master:} Facilitador del equipo, responsable de eliminar impedimentos (ej: Luis Huatay).\\
	extbf{Development Team:} Especialistas GIS, analistas técnicos, desarrolladores.

\subsection{Artefactos Scrum}
\subsubsection{(i) Product Backlog}
Lista priorizada de funcionalidades y requerimientos del proyecto, derivada del alcance inicial. Ejemplos de ítems:
\begin{itemize}
  \item Diseño de arquitectura de base de datos GIS.
  \item Desarrollo de mapas temáticos y dashboards.
  \item Implementación de integración con sistemas del cliente.
  \item Redacción de manuales técnicos y de usuario.
\end{itemize}

\subsubsection{(ii) Sprint Backlog}
Conjunto de tareas seleccionadas del Product Backlog para un sprint específico. Ejemplo para el Sprint 1:
\begin{itemize}
  \item T1: Inventariar sistemas del cliente para integración.
  \item T2: Diseñar prototipo de dashboard GIS.
  \item T3: Validar interoperabilidad con plataformas existentes.
\end{itemize}

\subsubsection{(iii) Increment}
Producto funcional entregado al final de cada sprint. Ejemplo: \\Sprint 1: Prototipo de dashboard GIS básico, con integración inicial a una fuente de datos.

\subsection{Descripción de eventos Scrum}
\subsubsection{(i) Sprint Planning}
	extbf{Objetivo:} Planificar el trabajo del sprint, seleccionando historias del Product Backlog.\\
	extbf{Duración:} 2 horas.\\
	extbf{Participantes:} Product Owner, Scrum Master, Equipo de Desarrollo.

\subsubsection{(ii) Daily (detalle de un Daily)}
	extbf{Fecha:} 12 de noviembre de 2025, 10:00 AM.\\
	extbf{Duración:} 15 minutos.\\
	extbf{Reportes:}
\begin{itemize}
  \item Ana: Ayer: Revisé requisitos de integración. Hoy: Diseñaré flujo de datos. Impedimentos: Ninguno.
  \item Carlos: Ayer: Avancé en desarrollo de API. Hoy: Terminaré pruebas unitarias. Impedimentos: Esperando acceso a servidor de pruebas.
\end{itemize}

\subsubsection{(iii) Review (cómo se llevará a cabo)}
	extbf{Objetivo:} Presentar el incremento al cliente y recibir feedback.\\
	extbf{Participantes:} Equipo Scrum, cliente, stakeholders clave.\\
	extbf{Actividades:} Demostración en vivo, discusión de mejoras, ajustes al backlog.

\subsubsection{(iv) Retrospective (de uno de los Sprint)}
	extbf{Sprint 1 – Técnica:} Start, Stop, Continue.\\
	extbf{Resultados:}
\begin{itemize}
  \item Start: Incluir pruebas de accesibilidad en el pipeline de QA.
  \item Stop: Implementar cambios sin aprobación previa del PO.
  \item Continue: Mantener Daily Scrums a las 10:00 AM.
\end{itemize}

\subsubsection{(v) Product Backlog Refinement}
	extbf{Objetivo:} Ajustar y priorizar ítems del backlog antes de cada sprint planning.\\
	extbf{Frecuencia:} Semanal.\\
	extbf{Responsable:} Product Owner con apoyo del equipo.

\subsection{Tablero Kanban del proyecto: Indicar todas las tareas y sus estados}


\section{III. Conclusiones}
El proyecto ``Consultoría Integral en Soluciones GIS'' representa una oportunidad estratégica para alinear los objetivos de Global GIS Service con las necesidades del mercado, mediante la implementación de una solución interoperable, escalable y centrada en el usuario. La aplicación del marco Scrum permitirá gestionar el desarrollo de manera iterativa e incremental, asegurando entregas tempranas, retroalimentación continua y adaptabilidad ante cambios. La combinación de una planificación robusta (PMI) con ejecución ágil (Scrum) garantizará no solo el cumplimiento de los plazos y presupuestos, sino también la satisfacción del cliente y el fortalecimiento del posicionamiento de GGS en el sector geoespacial latinoamericano.
\end{tabular}
\end{center}
\subsection{Logros Principales}
\begin{itemize}
    \item 100\% de mensajes pregrabados reemplazados en 15 cursos piloto.
    \item Transformación completa de 320 actividades de opcionales a obligatorias.
    \item Sistema de notificaciones implementado con 3 niveles de recordatorio.
    \item Dashboard docente con 5 reportes clave de seguimiento.
\end{itemize}

\section{Conclusiones y Lecciones Aprendidas}

\subsection{Lo que funcionó bien:}
\begin{itemize}
    \item Enfoque iterativo permitió ajustes basados en feedback real.
    \item Participación activa de usuarios finales en reviews.
    \item Comunicación transparente facilitada por SM.
\end{itemize}

\subsection{Áreas de mejora identificadas:}
\begin{itemize}
    \item Necesidad de mejores herramientas de prototipado rápido.
    \item Importancia de incluir soporte técnico desde planning.
    \item Valor de métricas cualitativas además de cuantitativas.
\end{itemize}

\subsection{Impacto Esperado:}
\begin{itemize}
    \item Mayor seriedad percibida en cursos virtuales.
    \item Incremento en tasa de completitud de actividades.
    \item Cultura de responsabilidad académica fortalecida.
\end{itemize}
